%قبل از شروع مسئله ترجیحاً چیزی ننویسید. اما اگر خواستید بنویسید، در قالب لیستی (فقرات و شمارش) بنویسید.
%برای اضافه کردن نمره، بعدی از شروع‌مسئله، [x] را اضافه کنید و به جای x  نمره مورد نظر را قرار دهید. برای زیرمسئله‌ها هم به همین صورت است.
%\شروع‌مسئله[10]
\شروع‌مسئله
% صورت اصلی مسئله را در این‌جا بنویسید.
به کمک قضیه اصلی توابع بازگشتی زیر را تحلیل کنید. در صورتی که تابعی با قضیه اصلی قابل تحلیل نباشد، فقط دلیل آن را ذکر کنید (نیازی به تحلیل این توابع نیست).

%اگر مسئله چند قسمت داشت به این صورت عمل کنید:
\شروع{زیرمسائل}
\زیرمسئله
$T(n) = 9T(\frac{n}{3}) + n^{2}$
\زیرمسئله
$T(n) = 3T(\frac{n}{5}) + n^{0.6}$
\زیرمسئله
$T(n) = 3T(\frac{n}{3}) + \frac{n}{\log{n}}$
\زیرمسئله
$T(n) = 16T(\frac{n}{4}) + n!$
\زیرمسئله
$T(n) = \frac{1}{4}T(\frac{n}{2}) + \frac{1}{n}$
\زیرمسئله
$T(n) = 2T(\sqrt[3]{n}) + (\log{n})^{2}$ \footnote{
راهنمایی : از تغییر متغیر $m = \log{n}$ استفاده کنید.
}
\پایان{زیرمسائل}

\solution{
	\شروع{زیرمسائل}
	\زیرمسئله 
	$
	\log_{3}{9} = 2
	$
	و داریم
	$f(n) = n^{2} = \Theta(n^2)$
	پس نتیجه میگیریم که 
	$T(n) = \Theta(n^2\log n)$
(حالت دوم از قضیه اصلی)
	\زیرمسئله $n^{\log_b a} = n^{0.682}$
	و چون 
	$n^{0.6}$
	از  
	$n^{0.682}$
	به صورت چند جمله ای  کوچکتر است پس در این حالت جواب 
	$\theta(n^{0.682})$
	است. (حالت اول قضیه اصلی)
	\زیرمسئله 
	چون که 
	$\log_{b}{a} = 1$
	و اینکه 
	$\frac{n}{\log{n}}$
	از $n^1$ به صورت چند جمله‌ای کوچکتر نیست، پس نمی‌توان از قضیه اصلی استفاده کرد.
	\زیرمسئله 
		$
	\log_{4}{16} = 2
	$
	و داریم
	$f(n) = n! = \Omega(n^{2+\epsilon})$
	پس نتیجه میگیریم که 
	$T(n) = \Theta(n!)$
	(حالت سوم از قضیه اصلی)
	\زیرمسئله 
	چونکه
	$\log_{2}{\frac{1}{4}} = -2$
	و اینکه 
	$\frac{1}{n}$
	از
	$\frac{1}{n^2}$
	به صورت چند جمله‌ای کوچکتر نیست، پس نمی‌توان از قضیه اصلی استفاده کرد.
	
	\زیرمسئله 
	 از تغییر متغیر
	 $m = \log{n}$
	  استفاده می‌کنیم و سپس قرار می‌دهیم
	   $T(2^m) = S(m)$.
	
	پس داریم:
	$S(m) = 2S(\frac{m}{3}) + m^{2}$
		در این‌جا از قضیه اصلی استفاده می‌کنیم و بدست میاوریم
		$S \in \Theta(m^2)$
		پس 
		$T \in \Theta(\log(m)^2$
	\پایان{زیرمسائل}
}





%%%%%%%%%%%%%%%%%%%%%%%%%%%%%%%%%%%%%%%%%%%%%%%%%%%%%
% برخی دستورات مفید برای تمرین‌ها. %
%%%%%%%%%%%%%%%%%%%%%%%%%%%%%%%%%%%%%%%%%%%%%%%%%%%%%

% فرمول‌های نمونه:
%$f(x) = \Theta(x)$
%$T(x, y) \leq \log(x) + 2^y + \binom{2n}{n}$
% {\tt code\_function}


% لیست موارد بدون شماره‌گذاری:
%\شروع{فقرات}
%    \فقره مورد اول 
%        \شروع{فقرات}
%            \فقره مورد اول داخلی 
%                \شروع{فقرات}
%                    \فقره اولین مورد خیلی داخلی 
%                    \فقره دومین مورد خیلی داخلی 
%                \پایان{فقرات}
%            \فقره مورد دوم داخلی 
%        \پایان{فقرات}
%    \فقره مورد دوم 
%\پایان{فقرات}

% لیست موارد شماره‌گذاری شده:
%\شروع{شمارش}
%    \فقره یک مورد 
%    \فقره یک مورد دیگر 
%    \فقره باز هم یک مورد دیگر
%\پایان{شمارش}

% برای نوشتن یک فرمول چند خطی که یک شماره بخوره. برای حالت‌های دیگر به سایت زیر مراجعه کنید:
% https://www.overleaf.com/learn/latex/Aligning_equations_with_amsmath
%\begin{align} 
%    \begin{split}
%        (x+y)^3 &= (x+y)^2(x+y) \\
%                &= (x^2+2xy+y^2)(x+y) \\
%                &= (x^3+2x^2y+xy^2) + (x^2y+2xy^2+y^3) \\
%                &= x^3+3x^2y+3xy^2+y^3
%    \end{split}                                 
%\end{align}

% نوشتن ماتریس و آرایه دوبعدی
%\begin{align}
%    A = 
%    \begin{bmatrix}
%        A_{11} & A_{21} \\
%        A_{21} & A_{22}
%    \end{bmatrix}
%\end{align}

% اضافه کردن تصویر با شرح. برای اضافه کردن تصویر بدون شرح، فقط دستور درج‌تصویر را استفاده کنید.
%\شروع{شکل}[h!]
%\تنظیم‌ازوسط
%  \درج‌تصویر[پهنا=0.5\پهنای‌متن]{img.jpg}
%\شرح{شرح تصویر}
%\برچسب{تصویر:۱}
%\پایان{شکل}
%\پایان{زیرمسائل}

%کد زیر را ببینید.
%%\lstinputlisting{source_filename.py}
%\begin{latin}
%\begin{lstlisting}
%# Import the modules
%import sys
%import random
%
%ans = True
%
%while ans:
%question = raw_input("Ask the magic 8 ball a question: (press enter to quit) ")
%
%answers = random.randint(1,8)
%
%if question == "":
%sys.exit()
%
%elif answers == 1:
%print "It is certain"
%
%elif answers == 2:
%print "Outlook good"
%
%elif answers == 3:
%print "You may rely on it"
%
%elif answers == 4:
%print "Ask again later"
%
%elif answers == 5:
%print "Concentrate and ask again"
%
%elif answers == 6:
%print "Reply hazy, try again"
%
%elif answers == 7:
%print "My reply is no"
%
%elif answers == 8:
%print "My sources say no"
%\end{lstlisting}
%\end{latin}
